% ============================================================
% Bibliography configuration
% ============================================================
% Three source files (see README.md):
%  - publications.bib      your papers that are part of the thesis (P1, P2, ...)
%  - publications_not.bib  your other papers during the PhD   (N1, N2, ...)
%  - bibliography.bib      everything else, cited entries only (1, 2, ...)
% The first two are read through generated/*_sorted.bib, which
% scripts/prepare_publication_bibs.py rebuilds on every latexmk run (it adds
% a year-month sortkey so they print oldest first). Never edit generated/.
\usepackage[
   backend=biber,
   style=numeric-comp,
   useprefix=true,
   sorting=none,
   sortcites=true,
   maxbibnames=99,
   maxnames=6,
   minnames=2,
   giveninits=true,
   defernumbers=true,
   backref=true
]{biblatex}

% Back-references ("Cited in pp. ...") closing each bibliography entry.
% The page numbers are hyperlinks, so they pick up linkcolor — the same
% green the cross-references use, which is what they are.
\DefineBibliographyStrings{english}{%
   backrefpage  = {Cited in p.},
   backrefpages = {Cited in pp.},
}
% Print it as its own short sentence rather than biblatex's default
% parenthetical. Entries already close on a parenthesised year, so the
% default "(2010). (Cited in p. 3)." stacks two paren groups and strands a
% period outside the second; here \finentry supplies the closing period.
\renewbibmacro*{pageref}{%
   \iflistundef{pageref}
      {}
      {\printtext{%
          \ifnumgreater{\value{pageref}}{1}
            {\bibstring{backrefpages}\ppspace}
            {\bibstring{backrefpage}\ppspace}%
          \printlist[pageref][-\value{listtotal}]{pageref}}}%
}

% Tag entries by datasource to split bibliography with stable precedence.
\DeclareSourcemap{
   \maps[datatype=bibtex,overwrite]{
      \map{
         \perdatasource{generated/publications_sorted.bib}
         \step[fieldset=usera, fieldvalue={pubsin}]
         \step[fieldset=presort, fieldvalue={a}]
      }
      \map{
         \perdatasource{generated/publications_not_sorted.bib}
         \step[fieldset=usera, fieldvalue={pubsnotin}]
         \step[fieldset=presort, fieldvalue={b}]
      }
      \map{
         \perdatasource{bibliography.bib}
         \step[fieldset=usera, fieldvalue={mainbib}]
         \step[fieldset=presort, fieldvalue={c}]
      }
      % Never truncate a co-first author into "et al.": an entry with an
      % N=cofirst keyword (see \thesis_bib_author_markers: below) shows at
      % least its first N names, overriding minbibnames=2 for entries past
      % maxnames=6 authors. The greedy .* captures the last (largest,
      % assuming ascending order) cofirst position. This replaces any
      % hand-set options field on such an entry.
      \map{
         \step[fieldsource=keywords, match=\regexp{\A.*(?:\A|,)\s*(\d+)\s*=\s*cofirst\s*(?:,|\Z)}, final]
         \step[fieldset=options, fieldvalue={minbibnames=$1}]
      }
      % Code packages, printed in their own "Code Packages" list (C1, C2,
      % ...) through the software driver below, which shows the repository
      % (stored in verba, since \AtEveryBibitem clears url) instead of a
      % journal code. To add a package: extend the citekey regex here and
      % add one map with its repository address (without https://).
      \map{
         \step[fieldsource=entrykey, match=\regexp{^(Doe:26a|Smith:26)$}, final]
         \step[fieldset=usera, fieldvalue={codepkg}]
         \step[typesource=article, typetarget=software]
         \step[typesource=misc, typetarget=software]
      }
      \map{
         \step[fieldsource=entrykey, match=\regexp{^Doe:26a$}, final]
         \step[fieldset=verba, fieldvalue={github.com/janedoe/MockPackage.jl}]
      }
      \map{
         \step[fieldsource=entrykey, match=\regexp{^Smith:26$}, final]
         \step[fieldset=verba, fieldvalue={github.com/asmith/OtherPackage}]
      }
      % Short journal names. Kept in sync with \bibjournaltitle below.
      \map{
         \step[fieldsource=journaltitle, match=\regexp{^Physical Review B$}, fieldset=shortjournal, fieldvalue={Phys. Rev. B}]
         \step[fieldsource=journaltitle, match=\regexp{^Physical Review Letters$}, fieldset=shortjournal, fieldvalue={Phys. Rev. Lett.}]
         \step[fieldsource=journaltitle, match=\regexp{^SciPost Physics$}, fieldset=shortjournal, fieldvalue={SciPost Phys.}]
      }
      % Link papers to their erratum via biblatex's related-entry mechanism
      % rather than \defbibentryset: entry sets never get a citeorder
      % position under sorting=none, so they always sort to the end of the
      % bibliography regardless of where they're actually first cited.
      % related keeps the main entry's own citeorder while still printing
      % the erratum folded into the same numbered item.
      \map{
         \step[fieldsource=entrykey, match=\regexp{^Smith:PRL20$}, final]
         \step[fieldset=related, fieldvalue={Smith:ErrPRL21}]
         \step[fieldset=relatedtype, fieldvalue={errata}]
      }
   }
}

% Filters by datasource tag (stored in usera).
\defbibcheck{pubsin}{\iffieldequalstr{usera}{pubsin}{}{\skipentry}}
\defbibcheck{pubsnotin}{\iffieldequalstr{usera}{pubsnotin}{}{\skipentry}}
\defbibcheck{mainbib}{\iffieldequalstr{usera}{mainbib}{}{\skipentry}}
\defbibcheck{codepkg}{\iffieldequalstr{usera}{codepkg}{}{\skipentry}}

% Separator printed before a related entry (e.g. an erratum) folded into
% the same numbered bibliography item.
\DeclareDelimFormat{begrelateddelim}{\addsemicolon\space}

% Oldest-first chronological sorting.
\DeclareSortingTemplate{oldestfirst}{%
   \sort[direction=ascending]{\field{sortkey}}
   \sort{\field{sorttitle}\field{title}}
   \sort{\field{sortname}\field{author}\field{editor}\field{translator}}
}

% Global sorting: P entries, then N entries, then everything else (presort
% a/b/c, set in the sourcemap above); the thesis's own papers go
% oldest-first by the sortkey that prepare_publication_bibs.py adds, the
% rest in citation order. sortcites orders the labels inside a \cite by
% this scheme, so a multi-key cite prints [P1, P3, 5] rather than in
% first-citation order. The P/N/main/C lists are printed separately with
% resetnumbers, so the main list's citation order is unaffected.
\DeclareSortingTemplate{thesis}{%
   \sort{\field{presort}}
   \sort[direction=ascending]{\field{sortkey}}
   \sort{\citeorder}
}
\ExecuteBibliographyOptions{sorting=thesis}

% Name punctuation: A. Author and B. Author; A. Author, B. Author, and C. Author
% (serial comma only for three or more names)
\DeclareDelimFormat{multinamedelim}{\addcomma\space}
\DeclareDelimFormat{finalnamedelim}{%
   \ifnumgreater{\value{liststop}}{2}{\addcomma}{}\space and\space}

% Bold thesis author in bibliography names (robust to accent variants in
% .bib files). Replace Doe/Jane/J. with your own family and given names,
% adding every spelling your .bib files use (e.g. Pay\'{a}, Payá, Paya).
\newcommand{\ifthesisauthorname}[2]{%
   \ifboolexpr{
      (test {\ifdefstring{\namepartfamily}{Doe}}) and
      (test {\ifdefstring{\namepartgiven}{Jane}} or
       test {\ifdefstring{\namepartgiven}{J.}} or
       test {\ifdefstring{\namepartgiven}{J}})
   }
      {#1}
      {#2}%
}
% Aux-file flag (same mechanism as \label/\ref) recording, per bib-file
% context (pubsin/pubsnotin), whether a cofirst-marked entry was actually
% printed — so the "Indicates co-first author" note above a list can be
% shown only when it applies to that list. Needs a rerun to pick up a
% newly-added cofirst entry, like any other aux-based cross-reference.
\makeatletter
\newif\ifthesiscofirstpubsin
\newif\ifthesiscofirstpubsnotin
% No "@" in this name (unlike the kernel macros in its body): it is called
% from inside the \ExplSyntaxOn block below, where "@" is not catcode-letter,
% so an "@"-bearing name would silently split into two tokens there.
\newcommand{\thesisMarkCofirstSeen}{%
   \@bsphack
   \iffieldequalstr{usera}{pubsin}
      {\global\thesiscofirstpubsintrue
       \protected@write\@auxout{}{\string\global\string\thesiscofirstpubsintrue}}
      {}%
   \iffieldequalstr{usera}{pubsnotin}
      {\global\thesiscofirstpubsnotintrue
       \protected@write\@auxout{}{\string\global\string\thesiscofirstpubsnotintrue}}
      {}%
   \@esphack
}
\newcommand{\thesisifcofirstseen}[2]{\csname ifthesiscofirst#1\endcsname #2\fi}
% The public \textsuperscript is broken by some TeX Live tagged-PDF kernel
% patches (it routes through an expl3 wrapper that never applies the
% size/raise, silently printing full-size inline text instead of a
% superscript). The kernel primitive \@textsuperscript sidesteps that, but
% positions the mark via a math-mode superscript, whose height follows
% fourier-otf's math axis/superscript parameters and sits noticeably higher
% than a typographic superscript should — so raise explicitly instead, at
% half the classic pre-2020 LaTeX kernel offset (.5ex), to bring it back down.
\newcommand{\thesisSuperscript}[1]{\leavevmode\raise.25ex\hbox{\fontsize\sf@size\sf@size\selectfont#1}}
\makeatother

% Author-position markers: keywords = {2=first} prints a * after author 2
% (first theory author), keywords = {2=cofirst,3=cofirst} a dagger after
% authors 2 and 3 (co-first authors).
\ExplSyntaxOn
\cs_new_protected:Npn \thesis_bib_author_markers:
 {
   \iffieldundef{keywords}
    {}
    {
      \tl_set:Nx \l_tmpa_tl { \thefield{keywords} }
      \int_set:Nn \l_tmpa_int { \value{listcount} }
      \exp_args:Nxx \regex_match:nnTF
         { (?:^|,)\s*\int_use:N \l_tmpa_int \s*=\s*first(?:\s*,|$) }
         { \tl_use:N \l_tmpa_tl }
         { \thesisSuperscript{*} }
         {}
      \exp_args:Nxx \regex_match:nnTF
         { (?:^|,)\s*\int_use:N \l_tmpa_int \s*=\s*cofirst(?:\s*,|$) }
         { \tl_use:N \l_tmpa_tl }
         { \thesisSuperscript{\textdagger}\thesisMarkCofirstSeen }
         {}
    }
 }
\newcommand{\bibauthormarkers}{\thesis_bib_author_markers:}
\ExplSyntaxOff
\renewcommand*{\mkbibnamegiven}[1]{\ifthesisauthorname{\mkbibbold{#1}}{#1}}
\renewcommand*{\mkbibnamefamily}[1]{\ifthesisauthorname{\mkbibbold{#1}}{#1}\bibauthormarkers}

% Keep bibliography compact.
% Books/chapters/proceedings keep url/isbn (needed to link the ISBN below);
% other entry types don't.
\AtEveryBibitem{%
   \ifboolexpr{
      test {\ifentrytype{book}} or
      test {\ifentrytype{incollection}} or
      test {\ifentrytype{inproceedings}}
   }
      {}
      {\clearfield{url}\clearfield{isbn}}%
   \clearfield{urldate}%
   \clearfield{month}%
   \clearfield{issn}%
   \clearlist{language}%
}

\newcommand{\bibtitlefont}[1]{{\itshape #1}}
% Fallback chain for short journal names; keep in sync with the sourcemap.
\newcommand{\bibjournaltitle}{%
   \iffieldundef{shortjournal}
      {\iffieldequalstr{journaltitle}{Physical Review B}
         {Phys. Rev. B}
         {\iffieldequalstr{journaltitle}{Physical Review Letters}
            {Phys. Rev. Lett.}
            {\iffieldequalstr{journaltitle}{SciPost Physics}
               {SciPost Phys.}
               {\printfield[noformat]{journaltitle}}}}}
      {\printfield[noformat]{shortjournal}}%
}
% \isdot marks a final abbreviation period (Phys. Rev. Lett.) as a dot, so
% biblatex does not swallow the comma before the volume.
\newcommand{\bibjournalblock}{%
   \printtext{\upshape\bibjournaltitle\isdot}%
   \setunit{\addcomma\space}%
   \printfield{volume}%
   \setunit{\addcomma\space}%
   \printfield[noformat]{pages}%
   \setunit{\space}%
   \printtext[parens]{\printfield{year}}%
}
\newcommand{\bibpreprintblock}{%
   \iffieldundef{number}
      {\printfield{eprint}}
      {\printfield{number}}%
   \setunit{\space}%
   \printtext[parens]{\printfield{year}}%
}
\newcommand{\bibpreprintlinkblock}{%
   \textcolor{my_red_d}{\bibpreprintblock}%
}

% ISBN, linked to the DOI if present, else to the URL, else left unlinked.
\newcommand{\bibisbnblock}{%
   \iffieldundef{isbn}
      {}
      {\iffieldundef{doi}
          {\iffieldundef{url}
              {\printfield{isbn}}
              {\href{\thefield{url}}{\printfield{isbn}}}}
          {\href{https://doi.org/\thefield{doi}}{\printfield{isbn}}}%
       \adddot}%
}

% Authors. Title. Edition, publisher, place, year. ISBN (linked) for books.
\DeclareBibliographyDriver{book}{%
   \usebibmacro{begentry}%
   \printnames{author}%
   \setunit{\adddot\space}%
   \printtext{\bibtitlefont{\thefield{title}}}%
   \adddot
   \setunit{\space}%
   \iffieldundef{edition}
      {}
      {\printfield{edition}\setunit{\addcomma\space}}%
   \printlist{publisher}%
   \setunit{\addcomma\space}%
   \printlist{location}%
   \setunit{\addcomma\space}%
   \printfield{year}%
   \adddot
   \setunit{\space}%
   \bibisbnblock%
   \setunit{\space}%
   \usebibmacro{pageref}%
   \usebibmacro{finentry}%
}

% Authors. Title. In Booktitle, ed. Editors, vol. N, pages. Publisher, place,
% year. ISBN (linked) for book chapters and conference proceedings.
\newcommand{\bibchapterentry}{%
   \usebibmacro{begentry}%
   \printnames{author}%
   \setunit{\adddot\space}%
   \printtext{\bibtitlefont{\thefield{title}}}%
   \adddot
   \setunit{\space}%
   \printtext{In\space}%
   \printfield[noformat]{booktitle}%
   \setunit{\addcomma\space}%
   \ifnameundef{editor}
      {}
      {\printtext{ed.\space}\printnames{editor}\setunit{\addcomma\space}}%
   \iffieldundef{number}
      {}
      {\printtext{vol.\space}\printfield{number}\setunit{\addcomma\space}}%
   \printfield[noformat]{pages}%
   \adddot
   \setunit{\space}%
   \printlist{publisher}%
   \setunit{\addcomma\space}%
   \printlist{location}%
   \setunit{\addcomma\space}%
   \printfield{year}%
   \adddot
   \setunit{\space}%
   \bibisbnblock%
   \setunit{\space}%
   \usebibmacro{pageref}%
   \usebibmacro{finentry}%
}
\DeclareBibliographyDriver{incollection}{\bibchapterentry}
\DeclareBibliographyDriver{inproceedings}{\bibchapterentry}

% Full journal code (journal, volume, pages, year) linked to DOI for articles.
\DeclareBibliographyDriver{article}{%
   \usebibmacro{begentry}%
   % Author-less entries (errata) skip the separator, which would otherwise
   % surface as a stray period after the title.
   \ifnameundef{author}
      {}
      {\printnames{author}%
       \setunit{\adddot\space}}%
   \printtext{\bibtitlefont{\thefield{title}}}%
   \printtext{\addcomma\space}%
   \iffieldundef{doi}
      {\bibjournalblock}
      {\href{https://doi.org/\thefield{doi}}{\bibjournalblock}}%
   % A related entry (an erratum) is printed before the closing period, so
   % the item reads "... (2020); Erratum ..., (2021)." rather than "(2020).;".
   \ifbibliography
      {\iftoggle{bbx:related}
         {\usebibmacro{related:init}%
          \usebibmacro{related}}
         {}}
      {}%
   \adddot
   \setunit{\space}%
   \usebibmacro{pageref}%
   \usebibmacro{finentry}%
}

% Code packages: author, title and repository, no journal code.
% The repository is stored without scheme and shown that way, linking to
% its https:// address.
\DeclareFieldFormat{codepkgurl}{\href{https://#1}{\nolinkurl{#1}}}
\DeclareBibliographyDriver{software}{%
   \usebibmacro{begentry}%
   \printnames{author}%
   \setunit{\adddot\space}%
   \printtext{\bibtitlefont{\thefield{title}}}%
   \setunit{\addcomma\space}%
   \printfield[codepkgurl]{verba}%
   \adddot
   \setunit{\space}%
   \usebibmacro{pageref}%
   \usebibmacro{finentry}%
}

% Full arXiv code (number + year) linked to DOI for preprints/misc items.
\DeclareBibliographyDriver{misc}{%
   \usebibmacro{begentry}%
   \printnames{author}%
   \setunit{\adddot\space}%
   \printtext{\bibtitlefont{\thefield{title}}}%
   \printtext{\addcomma\space}%
   \iffieldundef{doi}
      {\bibpreprintblock}
      {\href{https://doi.org/\thefield{doi}}{\bibpreprintlinkblock}}%
   \adddot
   \setunit{\space}%
   \usebibmacro{pageref}%
   \usebibmacro{finentry}%
}

\addbibresource{generated/publications_sorted.bib}
\addbibresource{generated/publications_not_sorted.bib}
\addbibresource{bibliography.bib}
